\chapter{Manufacturing Validation and Production Readiness}
\label{ch:manufacturing}

\textit{Proving the factory can reproduce and protect the qualified design.}

\textit{Read first:} production reference freeze; gauge R\&R; first-pass versus
final yield; process capability; ATP versus SPC; NPI gates.

\textit{Deep dive:} station correlation; escape classes; second-source evidence.

\textit{Reference:} NPI gate table; ATP control checklist; yield-owner split.

\begin{quote}
\noindent\textbf{Manufacturing validation at a glance}\\[0.4em]
\begin{itemize}
\item Freeze the production reference. Define the exact design, BOM, firmware,
      calibration, tools, and processes being validated.
\item Plan representative builds. Use production-intent materials, equipment,
      operators, suppliers, and stations.
\item Create unit genealogy. Preserve the material, process, software, and test
      history of every unit.
\item Validate the measurement system. Demonstrate accuracy, repeatability,
      reproducibility, stability, and station correlation.
\item Understand yield and distributions. Determine normal production behavior
      and identify dominant variation.
\item Validate ATP and screening. Prove that production tests detect relevant
      defects without excessive false rejects.
\item Establish process controls. Use SPC, audits, supplier controls, and
      reaction plans.
\item Ramp deliberately. Increase exposure only when evidence supports the next
      volume stage.
\item Close the loop. Feed escapes and process movement back into design,
      qualification, ATP, and supplier controls.
\end{itemize}
Manufacturing validation does not mean that one pilot lot passed. It means the
production system is defined, measurable, capable, traceable, controlled, and
prepared to react when it moves.
\end{quote}

A qualified engineering design is not automatically a manufacturable product.
Engineering units may be assembled by expert technicians, selected for favorable
components, calibrated individually, reworked repeatedly, and measured with
laboratory instruments that cannot support production takt time. A factory must
reproduce the same required behavior across materials, operators, tools, shifts,
sites, test stations, and time.

Manufacturing validation therefore begins by defining the production reference.
The team must know exactly which design, firmware, calibration, materials,
processes, fixtures, and test limits are being evaluated. It then builds
representative units using the intended process, preserves their genealogy, and
determines how much observed variation comes from the product versus the
measurement system.

Only after the measurements are trusted can yield, capability, and process
movement be interpreted. Production tests must then be shown to detect important
defects efficiently. Statistical process controls and reaction plans keep the
process stable as volume increases. Fleet and failure-analysis evidence complete
the loop by revealing escapes that factory evidence did not predict
(\Cref{ch:failure-modes,sec:fleet-triage}).

The aim is not to prove that every future unit will be good. The aim is to
establish a system that can repeatedly produce acceptable output, detect
unacceptable output, identify affected populations, and correct drift before it
becomes a fleet problem. The input is the design bounded by
\Cref{ch:reliability}.

\section{Manufacturing validation versus adjacent activities}
\label{sec:mfg-vs-adjacent}

\begin{table*}[ht]
\refstepcounter{table}\label{tab:mfg-vs-adjacent}%
\centering
\setlength{\tabcolsep}{3pt}%
\footnotesize
\renewcommand{\arraystretch}{1.2}%
\begin{tabular}{@{}p{0.16\linewidth}p{0.22\linewidth}p{0.28\linewidth}p{0.26\linewidth}@{}}
\toprule
Activity & Primary question & Typical evidence & Decision \\
\midrule
Design validation & Does the product meet intended system requirements? & Characterization, margin, interop & Approve system envelope \\
\addlinespace[0.8em]
Reliability qualification & Does the design survive named mechanisms and stresses? & Accelerated stress, degradation, confidence & Approve life/environment claim \\
\addlinespace[0.8em]
Manufacturing validation & Can production reproduce and measure the result? & Production-intent builds, MSA, yield, capability & Approve controlled ramp \\
\addlinespace[0.8em]
ATP & Can unacceptable units be detected economically? & Every-unit tests and validated proxies & Ship or reject unit \\
\addlinespace[0.8em]
SPC & Is the process remaining stable? & Time-ordered process metrics & Continue, contain, investigate \\
\addlinespace[0.8em]
Fleet monitoring & Does deployed behavior match the release model? & Telemetry, cohorts, returns & Expand, contain, improve \\
\bottomrule
\end{tabular}
\par\medskip
{\raggedright\footnotesize\textbf{Table~\thetable.} Adjacent activities. Do not
use qualification, manufacturing validation, ATP, SPC, and fleet monitoring
interchangeably. Manufacturing validation is Step~7 of the product-readiness
lifecycle in \Cref{sec:product-readiness}.\par}
\end{table*}
\FloatBarrier

\section{Supplier execution playbook}
\label{sec:supplier-exec}

The supplier path is milestones, performance targets, quality, and
manufacturability triage. That is not a soft skill. It is a concrete contract:
requirements, gates, acceptance tests, process control, and corrective action when
a lot goes wrong.

\whyexp{care about production lots?}{Because manufacturing escapes almost always
correlate with process history. Lot, date code, site, and firmware tags often
beat another night on one returned unit.}

\heuristic{Ask for the process change list before you invent new physics. Most
lot escapes sit next to a real change record.}

\tradeoff{Second source vs qualification burden}{Supply resilience and pricing
leverage}{Validation, interop matrix, and manufacturing differences}{When supply
or concentration risk exceeds the qual cost}{Qualify second sources based on
risk and evidence, not ideology.}

Lab hero samples versus manufacturable yield: same tradeoff as in
\Cref{ch:validation}. Optimize the system you can build, not the best bench unit.

\begin{dectree}
Design requirements
  |
Qualification
  |
ATP
  |
Production data
  |
Fleet data
  |
Failure analysis
  |
Updated limits or screens
  |
Next production cycle
\end{dectree}

Production validation is replayable and decision-oriented
(\Cref{sec:tree-production-loop}).

\paragraph{NPI as program maturity.}
New product introduction (\term{NPI}\firstterm{NPI}) is the cross-functional
program that moves a design into controlled production. It is an organizational
umbrella for engineering evidence, supplier readiness, factory process
development, test development, quality planning, operations, change control,
and ramp, not a technical evidence category beside the product-readiness
lifecycle in \Cref{sec:product-readiness}. Write exit criteria a supplier can
fail without ambiguity, not slogans. EVT/DVT/PVT/MP are program-phase labels,
not a calendar; dates and sample sizes belong in the program plan
(\Cref{sec:phase-labels}).

\paragraph{Why program maturity changes the question.}
Early in the program the question is whether the architecture can be made to
work at all (often during EVT), then whether behavior and margin are understood
across the intended corners (often during DVT). Once the part is characterized,
the question shifts to whether named failure mechanisms threaten life
(reliability qualification, Step~6), then whether production tooling and
suppliers can reproduce that qualified result across lots (manufacturing
validation, Step~7, often emphasized in PVT). A pilot asks whether laboratory
and factory assumptions survive a bounded, observable deployment. Mass
production is not another system-validation experiment: it is sustained control
with SPC, ECO, and RMA ownership
(\Cref{sec:ladder-mp,sec:ladder-fleet-monitor}).

Do not use an EVT hero sample as PVT evidence, and do not treat MP as fleet
monitoring alone. Hold a gate if the exit data are missing.
\Cref{tab:npi} retrieves the gate questions and release calls; it does not
define a second lifecycle beside the product-readiness Steps
(\Cref{sec:product-readiness,sec:tree-qualification}).

\begin{table*}[ht]
\refstepcounter{table}\label{tab:npi}%
\centering
\setlength{\tabcolsep}{4pt}%
\footnotesize
\renewcommand{\arraystretch}{1.25}%
\begin{tabular}{@{}p{0.10\linewidth}p{0.22\linewidth}p{0.36\linewidth}p{0.26\linewidth}@{}}
\toprule
Gate & Question & Representative evidence & Release decision \\
\midrule
EVT & Does it operate at all? & First light; CMIS bring-up; basic LIV/SMSR/RIN; one link closes BER & Continue / redesign integration \\
\addlinespace[1.25em]
DVT & Does it meet spec across corners? & Full ATP at $T$/$V$; prod-rep corners; stress plan + FIT model frozen & Enter qual / PVT / hold \\
\addlinespace[1.25em]
Qual & Env / reliability evidence ready? & Named mechanisms; sample plan; confidence (\Cref{sec:tree-qual-evidence}) & Enter PVT / hold \\
\addlinespace[1.25em]
PVT & Is it buildable at yield? & Multi-lot first-pass yield; MSA; ATP coverage; process capability; traceability; supplier readiness; validated rework; FAIR; production-host bring-up & Enter pilot / hold \\
\addlinespace[1.25em]
Pilot & Do assumptions hold in a bounded field trial? & Known serials/lots; enhanced telemetry; exit/rollback criteria & Open MP / restrict \\
\addlinespace[1.25em]
MP & Is quality sustained? & Steady DPPM; owned RMA Pareto; ECO control; fleet feedback & Keep shipping / CAPA / restrict \\
\bottomrule
\end{tabular}
\par\medskip
{\raggedright\footnotesize\textbf{Table~\thetable.} NPI gates (reference).
Decision unlocked: which release call the gate evidence supports. Pilot sits
between PVT and MP; MP is sustained control, not fleet monitoring alone.\par}
\end{table*}

\paragraph{Requirements and the production control plan are the contract.}
Requirements and the production control plan form the contract. Every
requirement needs a named evidence source and owner, but that evidence may be
ATP\firstterm{ATP}, a sampled audit, qualification, process control, supplier
evidence, or fleet monitoring. Version the requirements with the controls that
support them.

\heuristic{If a requirement has no evidence source, owner, or reaction plan, it
is a wish. ATP is only one possible control. If an ATP line has no requirement,
it is cost without a decision.}

\begin{enumerate}
\item \textbf{Requirements / PRD slice for the laser path:} fill
      \Cref{tab:laser-prd,sec:laser-reqs} (power class, grid, RIN@ORL, SMSR,
      derating, CMIS, FIT). Version it with the production control plan.
\item \textbf{Acceptance test plan (ATP):} the measurable tests that prove those
      requirements on every ship lot (or on a defined sample) when every-unit or
      sampled production test is the right control. Map each ATP line to a
      requirement; map life claims to qualification where appropriate
      (\Cref{sec:gr468}).
\end{enumerate}

\subsection{How production evidence is purchased}
\label{sec:atp-laser-read}

Every second of ATP costs line capacity. Every omitted screen leaves risk. Buy
evidence in categories that cannot substitute for each other.
\Cref{tab:atp-laser} is a working checklist for an EML pluggable or an ELSFP CW
module. Customize limits from the datasheet and the link budget; do not invent
numbers in the ATP itself.

\paragraph{Every-unit fast screens.}
Power class, CMIS state and bring-up, basic lane operation, and selected
electrical or optical checks that are cheap enough for 100\% coverage. These
catch dead units and firmware fails. They do not establish life.

\paragraph{Sampled expensive measurements.}
TDECQ, intrinsic and stressed RIN, detailed spectrum/SMSR, thermal corners, and
longer BER dwell. Use lot sample or audit when the signature is expensive and
correlated to escape risk.

\paragraph{Qualification-only stresses.}
HTOL, humidity, extensive cycling, and destructive analysis gather life or
mechanism evidence. They are not per-unit ship screens
(\Cref{sec:gr468,sec:tree-qual-evidence}).

\paragraph{Process controls.}
SPC on LIV, SMSR, RIN, TDECQ, and mate-cycle yield; golden-unit tracking; gauge
R\&R; incoming quality; first-article / FAIR after tooling or site change
(\Cref{sec:gauge-rr}).

\paragraph{Fleet controls.}
Telemetry, lot traceability, RMA codes split by mechanism, and recurrence
monitoring (\Cref{sec:fleet-triage}). These catch what ATP never saw.

A fast production power check does not establish life. HTOL does not prove every
shipped unit has correct firmware. SPC does not detect an unmeasured mechanism.
Source-level LIV, SMSR, and wavelength measurements can detect selected material
or device shifts when those measurements are directly available and correlated
to product risk. In a closed module, use a validated module-level proxy,
supplier evidence, sampled audit, or genealogy-based control; do not claim
internal measurement coverage when the production architecture does not expose
it. RIN and ORL protect the reflection environment; EAM/DCA protects Tx quality
on EML paths; CMIS protects field evidence. Burn-in may screen a demonstrated
infant-mortality population. ESD robustness is primarily established through
design qualification and handling controls; a finished-unit production screen
may not reliably detect latent ESD damage. Thermal class protects the derate
claim.

\begin{table*}[ht]
\refstepcounter{table}\label{tab:atp-laser}%
\centering
\setlength{\tabcolsep}{2.5pt}%
\footnotesize
\renewcommand{\arraystretch}{1.2}%
\begin{tabular}{@{}p{0.18\linewidth}p{0.16\linewidth}p{0.14\linewidth}p{0.22\linewidth}p{0.22\linewidth}@{}}
\toprule
Item & Method & Control class & Pass intent & Ties to \\
\midrule
LIV ($I_\mathrm{th}$, slope, kink) & SMU + power meter & 100\% ATP (source) / module proxy & kink-free bias window & wear-out, derate \\
\addlinespace[1.1em]
SMSR & OSA & 100\% ATP (source) / lot sample & single-mode vs.\ floor & modal noise \\
\addlinespace[1.1em]
Intrinsic RIN & PD + ESA & Lot sample / FA & quiet source floor & BER floor budget \\
\addlinespace[1.1em]
Stressed $\mathrm{RIN}_x\mathrm{OMA}$ & PD + ESA @ ORL & Lot sample / 100\% if escape & stressed Tx metric & named PMD \\
\addlinespace[1.1em]
Wavelength / grid & OSA / wavemeter & 100\% ATP or sample & channel ID & WDM lock \\
\addlinespace[1.1em]
Optical power class & power meter & 100\% ATP & class met & link budget \\
\addlinespace[1.1em]
EAM / TDECQ (EML) & bias + DCA & Lot sample / audit & ER, RLM, TDECQ & Tx quality \\
\addlinespace[1.1em]
CMIS / TWI bring-up & host / CMIS tool & 100\% ATP & state machine & telemetry \\
\addlinespace[1.1em]
Connector / ORL & mate + ORL meter & Periodic audit / sample & cycles + endface & packaging \\
\addlinespace[1.1em]
Burn-in (infant) & production screen & 100\% or lot sample & infant culled & not HTOL life \\
\addlinespace[1.1em]
HTOL life evidence & accelerated life & Qualification only & mechanism + FIT claim & GR-468 \\
\addlinespace[1.1em]
Driver/TIA ESD & JESD47 report & Qualification / audit & rating on file & IC reliability \\
\addlinespace[1.1em]
Thermal class & chamber & Lot sample / SPC & LIV/RIN/CMIS pass & derate \\
\addlinespace[1.1em]
Process monitors & SPC charts & SPC / process & drift detection & escape prevention \\
\addlinespace[1.1em]
Fleet cohort metrics & telemetry & Fleet telemetry & trend / alarm & recurrence \\
\bottomrule
\end{tabular}
\par\medskip
{\raggedright\footnotesize\textbf{Table~\thetable.} Control checklist for laser-bearing
modules (EML or ELSFP). Control class separates 100\% ATP, lot sample, audit,
qualification-only, SPC, and fleet telemetry. Intrinsic RIN and stressed
$\mathrm{RIN}_x\mathrm{OMA}$ are different metrics.\par}
\end{table*}
\FloatBarrier

\textbf{Exit for the ATP as a whole:} every ship lot (or defined sample) has
traceable pass data against versioned limits tied to the requirements slice.
\textbf{Decision:} ship, hold, or reopen DVT limits.

\paragraph{Incoming QC and SPC.}
Qual lots are small. Production catches drift that qual missed.

\begin{itemize}
\item \textbf{Incoming:} sample or 100\% screen against a subset of the ATP (at least
      power, CMIS, and a laser LIV/SMSR sample). Track DPPM by date code and site.
\item \textbf{SPC}\firstterm{SPC}: control charts on $I_\mathrm{th}$, slope, SMSR, Tx
      power, and burn-in fallout. A process shift is a hold, not a hope.
\item \textbf{First-article / FAIR:} when tooling, epi, or assembly site changes,
      rerun a defined FAIR package before open PO volume. Treat CMIS firmware revs
      the same way as process changes.
\end{itemize}

\paragraph{Excursions: 8D / CAPA.}
When a lot fails ATP, incoming, or field triage lands in the manufacturability
bucket (\Cref{sec:fleet-triage}), run structured corrective action:

\begin{enumerate}
\item \textbf{Contain:} quarantine WIP and ship holds; identify suspect date codes
      in the fleet.
\item \textbf{Evidence pack:} failing ATP rows, CMIS dumps, LIV/SMSR/RIN plots, DPA
      photos (facet, solder, FAU cross-section) compared to a golden unit.
\item \textbf{8D / CAPA}\firstterm{eightD}: confirmed mechanism with the supplier (process step,
      material lot, firmware), corrective action, and preventive control (ATP tighten,
      SPC limit, poke-yoke).
\item \textbf{Verify closure:} containment confirmed effective; mechanism reproduced
      or physically confirmed; corrective action removes the failure; no
      unacceptable regression introduced; production control detects recurrence;
      next lots remain stable; field cohort trend improves. Re-run FAIR alone is
      not enough for environmental, intermittent, or fleet-specific escapes
      (\Cref{sec:tree-evidence-block}).
\end{enumerate}

Do not close 8D on ``operator error'' without a control that would have caught it
at ATP or in process. If FA shows laser wear-out on a young unit, it may be a
reliability screen gap, not a supplier process bug; reclassify with
\Cref{sec:fleet-triage} before you argue FIT.

\paragraph{Milestone hygiene with partners.}
Align the partner calendar to gates, not slideware:

\begin{itemize}
\item freeze requirements before DVT samples are built;
\item freeze ATP limits before PVT yield is claimed;
\item freeze FIT/$E_a$ assumptions before reliability marketing numbers ship;
\item require ECO notice on laser die revision, TEC vendor, FAU epoxy, driver/TIA
      silicon revision (\Cref{sec:ic-reliability}), and CMIS firmware.
\end{itemize}

Your job in those meetings is to name the measurement that would kill the gate.
If nobody can point to an ATP row or a corner, the milestone is not real.

\keyidea{Manufacturing readiness is a frozen reference, a trusted measurement
system, multi-lot yield, ATP coverage, and owned SPC with reaction plans. Name
the mechanism family behind an escape
(\Cref{tab:qual-mechanism-families,sec:fleet-triage}) before you argue FIT or
open an 8D. Gate suppliers on ATP, multi-lot SPC, and FAIR. Do not run process
CAPA on wear-out, or FIT math on a dirty connector.}

\section{Yield analysis}
\label{sec:yield-analysis}

Yield is not one number. It splits by stage and by failure mode, and each split
points at a different owner.

\begin{description}
\item[Wafer / die yield] Process-limited: waveguide loss, ring resonance spread,
      heater shorts, photodiode dark current. Caught at wafer probe. Owner: foundry
      SPC.
\item[Assembly yield] Packaging-limited: fiber-array attach alignment, solder voids,
      wirebond pull, epoxy placement. Caught at module ATP. Owner: assembly supplier.
\item[Test yield (first-pass)] ATP-limited: units that fail one or more acceptance
      criteria on first pass. May include measurement-system false rejects (gauge
      RR). Owner: test engineering.
\item[Escaped DPPM (post-screen field failures)] Field or downstream failures that
      passed applicable production controls. Confirm a manufacturing escape only
      when evidence ties the mechanism to a preventable production or test-control
      gap; otherwise triage wear-out, interop, install, service, software, or
      residual latent risk. Owner: quality and reliability engineering.
\end{description}

\begin{table*}[ht]
\refstepcounter{table}\label{tab:yield-split}%
\centering
\setlength{\tabcolsep}{3.5pt}%
\footnotesize
\renewcommand{\arraystretch}{1.25}%
\begin{tabular}{@{}p{0.22\linewidth}p{0.28\linewidth}p{0.22\linewidth}p{0.22\linewidth}@{}}
\toprule
Yield stage & Main limit & First catch & Owner \\
\midrule
Wafer / die & Waveguide, resonance, heater, PD dark & Wafer probe & Foundry SPC \\
\addlinespace[1.5em]
Assembly & FAU align, solder, wirebond, epoxy & Module ATP & Assembly supplier \\
\addlinespace[1.5em]
Test (first-pass) & ATP fails; may include false rejects & ATP station + gauge RR & Test engineering \\
\addlinespace[1.5em]
Escaped DPPM & Passed screens; failed in fleet & Field RMA / triage & Quality / reliability \\
\bottomrule
\end{tabular}
\par\medskip
{\raggedright\footnotesize\textbf{Table~\thetable.} Yield stages, first catch, and
owner. Split escapes further in \Cref{tab:escape-classes}.\par}
\end{table*}
\FloatBarrier

Track yield by ATP row, lot, supplier site, tester, and date code. A yield drop
concentrated on one tester raises a measurement-system hypothesis. A yield drop
concentrated in one supplier lot raises a material or supplier-process
hypothesis, but station, shift, firmware, and chronology must be checked for
confounding before mechanism ownership is assigned. Operator correlation raises
a method, training, or equipment-interaction hypothesis. Date-code correlation
raises a change-history hypothesis. Correlation scopes the investigation; it is
not confirmed root cause. A yield drop with no observed correlation still needs
further investigation: verify gauge repeatability, expand stratification, and
test guardband or specification mismatch as hypotheses before concluding. Do not
open supplier corrective action until the measurement system is cleared
(\Cref{sec:hvm-test}).

\heuristic{Clear the tester with a golden unit before you escalate a supplier.
Station drift masquerades as a process excursion more often than engineers
admit.}

\usuallymeans{A golden unit fails on only one production station}{fixture,
calibration, cable, software limit, or operator path on that station}{a sudden
die-level failure of every good unit that station has ever seen}

\section{Process capability, limits, and guardbands}
\label{sec:process-capability}

Yield tells you how often units pass. Capability asks whether a stable process
can stay inside the requirement with room for measurement uncertainty and
expected drift. Before quoting capability, clear the measurement system
(\Cref{sec:gauge-rr}) and confirm the process is statistically stable enough for
the model.

Specification limits define acceptable product output. Control limits describe
the expected behavior of a statistically stable process. They are not
interchangeable: a stable process can be off-center and still produce bad units,
while an unstable process may temporarily remain inside specification. ATP
limits may be tighter than the customer specification as a guardband for
measurement uncertainty and expected drift, but the guardband must be justified
and must avoid double counting.

For a roughly normal, stable process,

\begin{equation}
C_p=\frac{\mathrm{USL}-\mathrm{LSL}}{6\sigma}
\end{equation}

\begin{equation}
C_{pk}=
\min\left(
\frac{\mathrm{USL}-\mu}{3\sigma},
\frac{\mu-\mathrm{LSL}}{3\sigma}
\right)
\end{equation}

$C_p$ measures potential spread if the process were centered. $C_{pk}$ also
reflects centering. Neither is meaningful until the process and measurement
system are stable enough for the model. Do not treat a single universal $C_{pk}$
threshold as automatic production readiness. Multimodal, drifting, censored, or
strongly non-normal optical data need more care than a textbook index.

\section{Production test at volume}
\label{sec:hvm-test}

\engcheck{Before production}{%
ATP $\cdot$
SPC $\cdot$
telemetry $\cdot$
supplier gates $\cdot$
monitoring owners $\cdot$
RMA-to-ATP feedback
(\Cref{sec:tree-checklists}).}

\subsection{Test time is a cost, coverage is a risk}

Every second in the acceptance test plan (ATP) times millions of units is line
capacity and real money. Every skipped measurement creates uncontrolled escape
risk; it is not automatically a field DPPM event (\Cref{sec:fit-dppm}). The core
tension in high-volume manufacturing is how much coverage you buy per second. The
expensive optical steps are thermal soak and corner runs, TDECQ on a sampling
scope, BER dwell long enough to trust a low pre-FEC target, laser burn-in, and
mate-cycle stress on ELSFP connectors. Some screens are statistical (sample
burn-in from a lot, audit TDECQ on a subset). Safety and enable-sequence faults
usually require 100\% coverage. At source or subassembly level, LIV, SMSR, and
wavelength may be economical direct controls when they are available and
correlated to product risk. In a closed module, use a validated module-level
proxy, supplier evidence, sampled audit, or genealogy-based control. Do not
claim internal measurement coverage when the production architecture does not
expose it (\Cref{sec:cmis,sec:laser-safety,tab:atp-laser}).

\tradeoff{More production screening vs cost}{Escape detection and earlier
catch}{Cycle time, tester cost, and false rejects that burn good units}{When a
named mechanism has a cheap, reliable detection signature}{Choose the cheapest
control that reliably detects the failure mode: 100\% ATP, sample audit, SPC,
or supplier process control.}

\tradeoff{Burn-in vs cycle time}{Infant-mortality removal when the screen
separates}{Line capacity, cost, and stress on healthy units}{When escape data
and mechanism justify the screen on this population}{Keep burn-in only while it
buys escapes you cannot catch cheaper elsewhere.}

\subsection{Where the test happens: wafer, die, module, system}

Push defect detection as far upstream as correlation allows. Wafer-level or PIC
probe catches process shifts (waveguide loss, ring resonance drift, bad heaters)
before fiber attach and packaging spend. Killing a bad die at probe is orders of
magnitude cheaper than an RMA (\Cref{sec:photonic-packaging}). Module ATP is the
full functional test: optical power class, TDECQ or proxy, sensitivity spot-check,
CMIS bring-up, and connector/ORL on ELS parts. System or golden-host bring-up
catches interop: media type, firmware rev, equalizer defaults, and the corners in
\Cref{sec:prod-corners}. Wafer test cannot catch fiber attach, FAU alignment,
epoxy creep, or connector wear. Those failures must survive to module ATP and,
for some signatures, to fleet telemetry (\Cref{sec:fleet-triage}).

\subsection{ATE-to-bench correlation and gauge R\&R}
\label{sec:gauge-rr}

Production testers are built for speed and cost, not lab fidelity. Correlation
asks whether ATE TDECQ, OMA, or sensitivity tracks the DCA/BERT reference within
a known offset and spread. Teach the measurement system explicitly:

\begin{description}
\item[Repeatability] Same station, operator, and unit.
\item[Reproducibility] Across stations, shifts, and operators.
\item[Bias] Production station versus reference lab.
\item[False reject / false accept] Guardband from measurement uncertainty.
\item[Golden-unit stability and station drift] Catch a stale golden or drifting
      ATE before it becomes a yield cliff or DPPM escape.
\end{description}

A golden unit is a station monitor, not a universal accuracy standard. It can
age, become contaminated, or be mishandled, so it needs controlled custody,
recertification, and retirement criteria. Use good, marginal, and failing units
across the measurement range for correlation, not one perfect center unit only.
Keep a golden module (and golden laser subassembly for ELS), run gauge R\&R
across testers and shifts, and correlate CMIS monitors to bench instruments the
same way you correlate TDECQ (\Cref{sec:cmis}). If the ATE and the DCA disagree,
fix the correlation before you argue with the supplier about spec.

\paragraph{Validating ATP coverage.}
Production-test coverage is validated using naturally failing units, controlled
parameter offsets, or carefully designed fault injection. The study should span
defect severity and measure detection probability, repeatability, false rejects,
station dependence, and test time. Fault injection validates only the
represented defect and severity range; it does not prove universal escape
coverage. Passing good units does not validate defect coverage.

\subsection{Screens, guardbanding, and SPC}

Burn-in (infant-mortality screen) and HTOL (life/mechanism evidence) trade
different risks against test time and cost
(\Cref{sec:wearout-modes,sec:gr468}). Test limits are usually guardbanded
tighter than the customer spec so field DPPM stays inside target under drift.
SPC control charts on LIV, SMSR, RIN, TDECQ, and mate-cycle yield by lot, site,
and date code catch a process shift before it becomes an 8D (\Cref{sec:supplier-exec}).
Production test is a yield, DPPM, and cost trade under a fixed reliability target.
It is not a pass/fail checkbox after the optics already work on a golden bench.

Keep the two rate languages with their owners. Detailed DPPM, yield-split, and
escape accounting belong to manufacturing and are developed in this chapter.
FIT and the life-rate arithmetic behind a reliability target belong to
qualification (\Cref{sec:fit-dppm}). A guardband is set by the first; the
target it protects comes from the second.


\section{Escaped defect analysis}
\label{sec:escaped-defects}

\begin{dectree}
Supplier escape containment flow
  |
Escape detected
  |
Provisional containment
  |
Scope analysis
  |
Refine contained population
  |
Investigate / confirm mechanism
  |
Corrective action
  |
Recurrence control
\end{dectree}

Provisional containment, scope refinement, mechanism confirmation, and
recurrence control are different actions. Apply
\Cref{sec:tree-escape,sec:tree-recurrence,sec:tree-evidence-block} to separate
immediate hold from FA and from the ATP, sample, SPC, or telemetry change that
closes the loop. Organize the spent margin with the five ledgers before naming
a component (\Cref{sec:laser-margin-erosion,sec:debug-fork}).

A potential production escape is a field or downstream failure that passed the
applicable production controls. It becomes a confirmed manufacturing escape only
when evidence connects the failure mechanism to a preventable production or
test-control gap. A field failure may instead be reliability wear-out, system
interoperability, installation damage, service condition, software, or residual
latent risk. When the failure is a confirmed preventable escape, post-screen
cases still split into two categories with different corrective actions:

\begin{table*}[ht]
\refstepcounter{table}\label{tab:escape-classes}%
\centering
\setlength{\tabcolsep}{3.5pt}%
\footnotesize
\renewcommand{\arraystretch}{1.25}%
\begin{tabular}{@{}p{0.20\linewidth}p{0.28\linewidth}p{0.24\linewidth}p{0.22\linewidth}@{}}
\toprule
Class & Meaning & Typical action & Lands in \\
\midrule
Preventable coverage & Screen or control could have caught it & Change recurrence control & Escape DPPM, CAPA \\
\addlinespace[1.5em]
Residual latent & No cost-effective screen & FIT / redundancy / replace & Residual FIT model \\
\bottomrule
\end{tabular}
\par\medskip
{\raggedright\footnotesize\textbf{Table~\thetable.} Escape classes. Preventable rows
change production; residual rows change the life model.\par}
\end{table*}
\FloatBarrier

\paragraph{Preventable coverage escapes.}
A defect that a cost-effective screen or process control could have caught but
did not, because the control was absent, miscalibrated, or insufficiently
stressed. A preventable escape requires a changed recurrence control. The best
control may be upstream process control, supplier control, design poka-yoke,
incoming inspection, sampled audit, ATP, SPC, or service procedure: the earliest
reliable and economical point of prevention or detection. Common mechanisms in
optical modules:

\begin{itemize}
\item \textbf{Contamination:} particles trapped during assembly that only move under
      thermal cycling or vibration, shifting coupling or raising ORL intermittently.
\item \textbf{Marginal alignment:} fiber-attach or FAU position within spec at room
      temperature but outside at the thermal corner, creating a
      temperature-dependent loss that the ATP chamber did not run.
\item \textbf{Weak solder or wirebond:} passes pull test but fatigues under
      repeated thermal cycling, creating intermittent opens or increased resistance.
\item \textbf{Firmware corner:} a state-machine path or calibration table entry that
      production never exercises but the fleet hits on a specific boot sequence,
      temperature bin, or host interaction.
\item \textbf{Packaging stress:} residual mechanical stress from underfill cure or
      lid attach that relaxes over time, shifting alignment or birefringence.
\end{itemize}

For each preventable escape, trace the failure signature back to the earliest
reliable control point in the production flow and change the recurrence control
there. Do not assume the answer is always a new or tighter finished-unit ATP
line.

\heuristic{An escape without a changed recurrence control is unfinished work.
Containment stops the bleed; the control stops the next lot.}

\paragraph{Residual latent failures.}
A defect that no cost-effective screen can separate from good units at the time of
test. Examples include cosmic-ray-induced single-event upsets, rare material
inclusions below detection limits, and process-marginal units that sit inside all
guardbands but fail under a specific field combination of temperature, ORL, and
neighbor load. These go into the residual FIT model, not into ATP. Document why no
reasonable screen exists so the decision is auditable and revisitable as test
technology improves.


\section{Worked case study: yield loss blamed on the laser supplier}
\label{sec:mfg-case-yield}

\textit{Illustrative numbers only.} A 240-unit production-intent build shows 90\%
first-pass yield. Low OMA and high laser bias dominate the Pareto. Failures
appear to cluster on one laser date code. Two ATP stations disagree by about
$0.4$~dB on the same units. Retest recovers many fails.

\begin{enumerate}
\item \textbf{Verify measurement first.} Golden units and station-to-station
      correlation show station~B reads low by $\sim0.4$~dB. Do not open laser
      supplier CAPA yet (\Cref{sec:gauge-rr}).
\item \textbf{Stratify the population.} Split by station, fixture, operator,
      shift, and laser lot. Lot and station are confounded: the suspect date code
      ran mostly on station~B.
\item \textbf{Controlled swap.} Move the same units and fixture across stations.
      The OMA offset follows the station, not the laser lot.
\item \textbf{Find the mechanism.} Fixture insertion loss has drifted. Repair and
      recertify the station; remeasure the held population.
\item \textbf{Quantify residual supplier difference.} After station repair, one
      laser lot still sits slightly high in bias. Contain that lot, tighten
      incoming LIV sample, and document the residual as a material signal rather
      than the original yield cliff.
\item \textbf{Update controls.} Golden-unit cadence, gauge R\&R, and SPC on
      station offset resume. Restricted ramp continues only after a confirmation
      lot clears first-pass yield and station agreement.
\end{enumerate}

A yield problem can belong to the product, material, process, measurement
system, software, or data. Verify measurement before redesigning the product or
blaming the supplier.


\section{Engineering lens}
\label{sec:manufacturing-engineering-lens}

\subsection{How it works}

Manufacturing validation proves that the qualified design
(\Cref{ch:reliability}) can be reproduced and protected at scale: frozen
reference, representative builds, genealogy, trusted measurement, yield and ATP
coverage, SPC with reaction plans, and staged ramp.

\subsection{How it is measured}
\label{sec:manufacturing-how-measured}

Production records first-pass and final yield, retest and rework rates, fallout
by ATP row, measurement distributions, gauge repeatability, station correlation,
and escaped defects per million. Keep the chain from incoming material through
module ATP so a drift can be traced to its first observable point
(\Cref{sec:hvm-test,sec:supplier-exec,sec:fleet-triage}).

\subsection{How it fails}
\label{sec:manufacturing-how-fails}

Programs fail at scale through variation and escapes. Variation produces a weak
tail across wafer, lot, site, or assembly line. An escape is a defect the current
screen cannot see or a control that was not run. Yield can drop because the
product changed, the process moved, incoming material moved, calibration
changed, or the tester moved. Do not open supplier corrective action until the
measurement system is cleared.

\failuremode{Yield drop}{First-pass yield falls from its stable baseline, often on
one ATP row, tester, lot, or shift.}{Process drift, supplier-lot variation,
assembly change, calibration or fixture drift, software revision, or a changed
guardband.}{Pareto by test and lot, golden-unit history, gauge repeatability,
station correlation, incoming data, and destructive analysis on selected
failures.}{Contain suspect material, clear the tester, identify the first changed
input, correct it, and verify with a controlled lot before release.}

\subsection{How it is debugged}
\label{sec:manufacturing-how-debugged}

For a yield fall, freeze software, limits, and suspect material. Split by tester,
shift, lot, supplier site, and ATP row. Golden unit across stations; failing unit
on a reference bench. Station-follows means repair the measurement system;
unit-follows means upstream process and mechanism FA. Contain first, confirm
second, change the process third, verify on fresh data.

\debugstory{Module yield fell sharply after a supplier lot change.}{The failure
Pareto pointed to one-lane TDECQ. Golden units passed all stations, and failed
units kept the bad lane on the reference bench. Cross-sections showed a shifted
fiber-array attach.}{The electrical path and testers were stable.}{An assembly
fixture change moved one fiber row outside its coupling window.}{The lot was held,
the fixture was restored, first-article coupling checks were tightened, and the
supplier control plan was revised.}

\section{Interview takeaway}
\label{sec:manufacturing-design-review}

\keyidea{Manufacturing validation proves that a qualified design can be
reproduced and protected at scale. Freeze the production reference, build
representative lots, preserve genealogy, validate the measurement system,
understand distributions and first-pass yield, prove ATP coverage, and establish
SPC with reaction plans. Increase volume only as evidence supports greater
exposure. The goal is not one successful lot; it is a production system that
remains capable, traceable, measurable, and correctable.}

Junior mistake: escalate a supplier before the measurement system is cleared, or
treat two hand-selected lots as multi-lot evidence
(\Cref{sec:gauge-rr,ch:failure-modes,app:case-studies}).

\subsection{Interview Q\&A: Manufacturing Validation}
\label{sec:manufacturing-interview-qa}

Practice speaking these answers aloud. Prefer first-person reasoning over tool
lists. Detail lives in
\Cref{sec:yield-analysis,sec:process-capability,sec:gauge-rr,sec:hvm-test,tab:mfg-vs-adjacent}.

\paragraph{Question 1. What is manufacturing validation, and how does it differ
from reliability qualification, ATP, SPC, and fleet monitoring?}
\textit{Tests:} terminology, lifecycle ownership, and decision boundaries.

\textit{Spoken answer.}
``Manufacturing validation asks whether the intended production system can
repeatedly reproduce, measure, detect, trace, and control the qualified design.
Reliability qualification asks whether representative hardware survives named
lifetime and environmental mechanisms. ATP makes a unit-level ship or reject
decision using validated production measurements or proxies. SPC asks whether
the process is moving over time, even before units fail specification. Fleet
monitoring checks whether deployed populations match the release model. They
share data, but they answer different questions and unlock different decisions''
(\Cref{ch:reliability,ch:validation,tab:mfg-vs-adjacent}).

\textit{Pressure follow-up.} ``Can a product pass qualification and still fail
manufacturing validation?''\\
\textit{Answer pivot.} ``Yes. Hand-built qualification samples may be reliable
while the production line has weak alignment yield, poor test correlation,
uncontrolled rework, or supplier-lot variation. Qualification evidence does not
establish factory repeatability.''

\textit{Trap:} ``Manufacturing validation is qualification performed on
production units.''

\paragraph{Question 2. What must be frozen before a production-intent validation
build?}
\textit{Tests:} configuration control and interpretability.

\textit{Spoken answer.}
``I would freeze enough of the production reference that the build has an
interpretable configuration. That includes the hardware and BOM revisions,
approved suppliers, firmware, CMIS behavior, calibration algorithm,
manufacturing recipes, tooling, work instructions, fixtures, test software,
reference planes, and acceptance limits. Deviations can exist, but they must be
explicit and traceable. If design, firmware, calibration, and process are all
changing simultaneously, I may still run an engineering learning build, but I
would not call its yield production-validation evidence.''

\textit{Pressure follow-up.} ``Does design freeze mean no changes are
allowed?''\\
\textit{Answer pivot.} ``No. It means changes are controlled, versioned, and
tied to the units affected. The purpose is causality and traceability, not
bureaucratic immobility.''

\textit{Trap:} ``The schematic revision and BOM are frozen, so the factory is
ready.''

\paragraph{Question 3. How would you choose the size and structure of a
manufacturing-validation build?}
\textit{Tests:} representative sampling, confounding, and release evidence.

\textit{Spoken answer.}
``I would choose the build around the sources of variation and the decision it
must support, not around one universal sample count. I want representative
production lots, component date codes, suppliers, operators, shifts, tools,
fixtures, stations, and product variants. A few hundred units may be useful for
early distributions and line behavior, but that number does not establish a very
low escape rate. I would prefer a balanced build that separates variables over a
larger build where one supplier lot always runs on one station and one shift.''

\textit{Pressure follow-up.} ``Would 200 units be enough?''\\
\textit{Answer pivot.} ``Enough for what? Two hundred may characterize early
yield and station behavior, but the confidence depends on lot diversity, defect
frequency, measurement capability, and the release decision. I would state what
those 200 units do and do not establish.''

\textit{Trap:} ``We normally build 200 units because that is statistically
significant.''

\paragraph{Question 4. What traceability and unit genealogy do you need?}
\textit{Tests:} population scoping and evidence preservation.

\textit{Spoken answer.}
``I want each serial number connected to its product revision, firmware,
calibration version, material lots, supplier sites, process tools, operator or
station where relevant, timestamp, test-software revision, fixture, raw
measurements, applied limits, and complete retest or rework history. A timestamp
alone is not genealogy. I also preserve the first failure rather than
overwriting it with the eventual passing result. That allows a field or yield
issue to be scoped by lot, tool, station, firmware, or process change instead of
treating every failure as an isolated unit.''

\textit{Pressure follow-up.} ``Why preserve every retest if the unit finally
passes?''\\
\textit{Answer pivot.} ``Because repeated retest may reveal measurement
instability, intermittent product behavior, marginality, or an undocumented
intervention. Final pass alone destroys the evidence needed to improve the
process.''

\textit{Trap:} ``A serial number, final ATP result, and build date are
sufficient.''

\paragraph{Question 5. What does measurement-system analysis tell you, and how
would you validate a production station?}
\textit{Tests:} accuracy, repeatability, reproducibility, correlation, and
golden-unit limits.

\textit{Spoken answer.}
``Before interpreting process variation, I need to know how much observed
variation comes from the measurement system. I would evaluate repeatability on
the same station and unit, reproducibility across stations, operators, or
fixtures, bias relative to a trusted laboratory reference, stability over time,
and error across the measurement range. For optical ATP I would correlate
representative good, marginal, and failing units between the production station
and the DCA, BERT, OSA, or power-meter reference. Golden units help detect
drift, but they also need custody, recertification, and retirement criteria''
(\Cref{sec:gauge-rr}).

\textit{Pressure follow-up.} ``What does gauge R\&R not tell you?''\\
\textit{Answer pivot.} ``It does not prove the product meets its specification
or lifetime requirement. It tells me whether the measurement system can resolve
the variation needed for the production decision.''

\textit{Trap:} ``The station passed calibration, so its production measurements
are trustworthy.''

\paragraph{Question 6. Explain first-pass yield, final yield, retest, rework,
and why distributions matter.}
\textit{Tests:} yield integrity and hidden process instability.

\textit{Spoken answer.}
``First-pass yield is the fraction that passes without retest or product
intervention. Final yield includes units recovered through valid retest or
approved rework. I report both, along with retest, rework, scrap, and
invalid-test rates. A high final yield can hide a weak process if many units
fail initially or require repeated adjustment. I also examine the parameter
distributions, because two lines can have the same yield while one has a
centered narrow distribution and the other has a wide tail clipped by the
acceptance limit''
(\Cref{sec:yield-analysis}).

\textit{Pressure follow-up.} ``A line has 99\% final yield and 85\% first-pass
yield. Is it healthy?''\\
\textit{Answer pivot.} ``Not without explanation. The recovery path may be
tester instability, uncontrolled tuning, or product marginality. I would
investigate the first-pass Pareto and retest behavior before approving the
process.''

\textit{Trap:} ``Only final yield matters because those are the units that
ship.''

\paragraph{Question 7. How would you design the production-test architecture?}
\textit{Tests:} every-unit screens, sampled audits, process controls, and test
economics.

\textit{Spoken answer.}
``I would start with the defect or risk, then choose the earliest and least
expensive control that detects or prevents it reliably. Every-unit ATP should
cover fast, high-value checks such as identity, firmware, CMIS states, basic
optical and electrical function, power, wavelength, alarms, and selected BER or
quality proxies. Expensive measurements such as full temperature sweeps, long
BER waterfalls, detailed TDECQ, RIN, ORL sensitivity, or destructive inspection
may belong in sampled audits. Some mechanisms have no useful finished-unit
screen and must be controlled by design, qualification, supplier controls, or
process monitoring''
(\Cref{sec:hvm-test}).

\textit{Pressure follow-up.} ``How do you prove an ATP screen actually catches
the defect?''\\
\textit{Answer pivot.} ``I use naturally failing units or controlled known-defect
injection across the relevant severity range, then measure detection
probability, false rejects, station dependence, and test time. Passing good
units does not validate defect coverage.''

\textit{Trap:} ``I would put every engineering test into ATP so nothing
escapes.''

\paragraph{Question 8. Explain specification limits, control limits,
capability, and guardbands.}
\textit{Tests:} statistical process interpretation and limit discipline.

\textit{Spoken answer.}
``Specification limits define acceptable product output. Control limits describe
the expected behavior of a statistically stable process. They are not
interchangeable: a stable process can be off-center and produce bad units, while
an unstable process may temporarily remain inside specification. $C_p$ describes
potential spread relative to the specification, while $C_{pk}$ also reflects
centering, but neither is meaningful until the process and measurement system
are stable enough for the model. ATP guardbands may be tighter than the customer
limit to cover measurement uncertainty and expected drift, but the guardband must
be justified and must avoid double counting''
(\Cref{sec:process-capability}).

\textit{Pressure follow-up.} ``Can you relax an ATP limit because yield is
poor?''\\
\textit{Answer pivot.} ``Only after re-establishing the link between the
customer requirement, measurement uncertainty, process distribution, and escape
risk. Yield recovery alone is not a technical justification.''

\textit{Trap:} ``If $C_{pk}$ is greater than 1.33, the process is automatically
ready for production.''

\paragraph{Question 9. Yield drops suddenly. Walk me through your response.}
\textit{Tests:} containment, tester clearing, stratification, and controlled
confirmation.

\textit{Spoken answer.}
``I would preserve the timeline and immediately bound exposure while avoiding
premature root-cause claims. I would freeze test software, limits, firmware, and
suspect material, then clear the measurement system using golden units and
reference-bench correlation. Next I would stratify first-pass failures and
parameter distributions by ATP row, station, fixture, operator, shift, material
lot, supplier site, firmware, and build order. Correlation creates a leading
hypothesis, not confirmation. I would run the smallest controlled swap or
experiment that separates tester, product, material, and process ownership, then
correct and verify on fresh production data before releasing the held
population''
(\Cref{sec:mfg-case-yield,sec:escaped-defects}).

\textit{Pressure follow-up.} ``The failures correlate strongly with one laser
lot. Do you open supplier CAPA?''\\
\textit{Answer pivot.} ``I may notify and provisionally contain the lot, but I
would first check whether the lot is confounded with one station, tool, or
shift. I escalate the supplier with evidence, not merely a Pareto correlation.''

\textit{Trap:} ``The lot correlation proves the laser supplier caused the yield
loss.''

\paragraph{Question 10. What is the difference between ATP and SPC, and what
makes an SPC program useful?}
\textit{Tests:} unit decisions versus process-time decisions.

\textit{Spoken answer.}
``ATP decides whether an individual unit meets the production acceptance
criteria. SPC monitors selected process or product metrics in build order to
detect movement before it becomes an escape or yield cliff. A useful SPC metric
has a trustworthy measurement, sensitivity to a real process input, an owner, a
trigger, an immediate containment window, an investigation path, and restart
criteria. A control chart without a reaction plan is just a visualization.''

\textit{Pressure follow-up.} ``The process is inside specification but shows a
sustained upward trend in laser bias. What do you do?''\\
\textit{Answer pivot.} ``I treat the trend as evidence of process movement even
before units fail specification. I contain according to the reaction plan, verify
the station, and investigate materials, calibration, temperature, and process
chronology.''

\textit{Trap:} ``SPC means rejecting any unit outside the product
specification.''

\paragraph{Question 11. How would you qualify a second source or manage a
supplier change?}
\textit{Tests:} supplier evidence, equivalence, change control, and
qualification re-entry.

\textit{Spoken answer.}
``I first define what equivalence means for the change. A second-source
component may affect performance distributions, calibration, thermal behavior,
reliability mechanisms, assembly interaction, and ATP correlation. A
second-source module adds firmware, CMIS, interoperability, telemetry, and
manufacturing-system differences. I would compare representative lots, margins,
process capability, measurement correlation, qualification evidence, and
supported operating corners. Any supplier, site, material, firmware, or tooling
change should have traceable pre- and post-change populations and a defined
revalidation or requalification plan''
(\Cref{ch:reliability,sec:supplier-exec}).

\textit{Pressure follow-up.} ``The supplier says the replacement is
form-fit-function equivalent. Is that enough?''\\
\textit{Answer pivot.} ``No. Form, fit, and nominal function do not establish
distribution, margin, reliability mechanism, process interaction, or field
behavior. The evidence depth should match what the change can affect.''

\textit{Trap:} ``Once the supplier is approved, future process changes are their
responsibility.''

\paragraph{Question 12. Give me a 60-second manufacturing-validation plan for a
new optical module.}
\textit{Tests:} complete Staff-level manufacturing answer.

\textit{Spoken answer.}
``I would begin by freezing the production reference: design, BOM, suppliers,
firmware, calibration, manufacturing recipes, fixtures, test software, and
limits. Then I would plan balanced production-intent builds that cover relevant
lots, operators, tools, stations, and material variation, with complete unit
genealogy. Before interpreting yield, I would validate production measurements
through gauge R\&R, station correlation, and golden-unit stability. I would
analyze first-pass yield and parameter distributions, validate ATP coverage with
known-defect units or controlled fault injection, and establish SPC metrics with
owners and reaction plans. Volume would ramp in controlled stages, with supplier
change control, escape containment, and fleet feedback tied back to ATP,
process, qualification, or design.''

\textit{Pressure follow-up.} ``What evidence would make you hold the ramp?''\\
\textit{Answer pivot.} ``I would hold for unresolved measurement disagreement,
missing genealogy, unexplained yield or parameter movement, weak ATP coverage
for a high-impact defect, uncontrolled rework, supplier changes without
evidence, or field behavior inconsistent with the release model.''

\textit{Trap:} ``I would build a pilot lot, confirm acceptable yield, and open
mass production.''

Score each response using the shared chapter-interview rubric in
\Cref{sec:chapter-interview-rubric}. Repeat any answer that does not identify
the production decision, the evidence required, and the reaction if the evidence
fails.
